\documentclass[11pt]{article}
\usepackage{fontspec}
\setmainfont{Latin Modern Roman}
\usepackage{amsmath,amssymb}
\usepackage{booktabs}
\usepackage[margin=2.5cm]{geometry}
\usepackage{hyperref}

\title{MDraft LaTeX performance: small}
\author{MDraft benchmarks}
\date{}

\begin{document}
\maketitle

\section{A short section}
\label{sec:small}

Performance probe paragraph S1. This single-file document is the small case in
MDraft's LuaLaTeX benchmark. It contains prose, one equation and one table, so
a successful compile proves the engine rendered real content rather than an
empty page.

\begin{equation}
  \label{eq:small:sum}
  \sum_{k=1}^{n} k = \frac{n(n+1)}{2}.
\end{equation}

Equation~\ref{eq:small:sum} is valid for every positive integer $n$. The same
identity appears in many introductions to induction, and it is short enough to
check by hand for small values.

\begin{table}[h]
  \centering
  \begin{tabular}{lrr}
    \toprule
    Item & Count & Cost \\
    \midrule
    Alpha & 3 & 12.50 \\
    Beta  & 5 & 21.00 \\
    Gamma & 8 & 33.75 \\
    \bottomrule
  \end{tabular}
  \caption{A small table with valid content.}
  \label{tab:small:costs}
\end{table}

Table~\ref{tab:small:costs} is rendered from the same source. The totals are
not used by the benchmark; the table exists so the compile has a real float.

\section{A second short section}
Performance probe paragraph S2. The benchmark edits a letter in this paragraph
and waits for the preview to commit the new PDF. Nothing here depends on the
machine's fonts beyond Latin Modern, so the result is reproducible.

A second paragraph gives the editor a stable offset and keeps the document from
being a single screen. The section is intentionally short: the small fixture
measures the engine's fixed cost, not a long typeset.
\end{document}
